%%%%%%%%%%%%%%%%%%%%%%%%%%%%%%%%%%%%%%%%%%%%%%%%%%%%%%%%%%%%%%%%%
% GIÁO ÁN STEM GIAI ĐOẠN 2: SMARTLOGI - TỐI ƯU HÓA VẬN TẢI
% Phiên bản đầy đủ - Tương thích với stem.tex Giai đoạn 1
%%%%%%%%%%%%%%%%%%%%%%%%%%%%%%%%%%%%%%%%%%%%%%%%%%%%%%%%%%%%%%%%%
\documentclass[12pt,a4paper]{article}

% --- 1. GÓI CƠ BẢN ---
\usepackage[utf8]{inputenc}
\usepackage[T1]{fontenc}
\usepackage[vietnamese]{babel}
\usepackage{graphicx}
\usepackage{amsmath,amssymb}
\usepackage{array}
\usepackage{booktabs}
\usepackage{multirow}
\usepackage{tabularx}
\usepackage{wrapfig}
\usepackage{ragged2e}
\usepackage{setspace}
\usepackage{tikz}
\usetikzlibrary{shapes,arrows,positioning,patterns,calc}
\usepackage{fontawesome5}
\usepackage{verbatim}
\usepackage{colortbl}

% --- 2. FONT ---
\usepackage{newtxtext, newtxmath}
\usepackage[scaled=0.92]{helvet}
\usepackage{courier}
\usepackage{microtype}

% --- 3. MÀU SẮC ---
\usepackage[svgnames]{xcolor}
\definecolor{primary}{RGB}{0,102,204}
\definecolor{secondary}{RGB}{80,80,80}
\definecolor{accent}{RGB}{255,69,0}
\definecolor{success}{RGB}{34,139,34}
\definecolor{warning}{RGB}{255,140,0}
\definecolor{TableHeader}{gray}{0.9}

\usepackage[left=2cm,right=2cm,top=2.5cm,bottom=2cm,headheight=15.6pt]{geometry}
\usepackage{fancyhdr}
\pagestyle{fancy}
\fancyhf{}
\renewcommand{\headrule}{\color{primary}\hrule width\textwidth height 1.5pt}
\fancyhead[L]{\color{primary}\footnotesize\sffamily\textbf{SỞ GD\&ĐT TP.HCM - HỘI THI STEM/STEAM CẤP THÀNH PHỐ}}
\fancyhead[R]{\color{primary}\footnotesize\sffamily\textbf{NĂM HỌC 2025-2026}}
\fancyfoot[C]{\sffamily\thepage}
\renewcommand{\footrule}{\color{primary}\hrule width\textwidth height 0.8pt}

% --- 4. GÓI HỖ TRỢ ---
\usepackage{enumitem}
\usepackage{tcolorbox}
\tcbuselibrary{skins, breakable}
\usepackage[colorlinks=true,linkcolor=primary,urlcolor=accent,citecolor=success,
    pdftitle={Giao an STEM: SmartLogi Logistics},pdfauthor={Nguyen Van Sang}]{hyperref}

% --- 5. TIÊU ĐỀ ---
\usepackage{titlesec}
\titleformat{\subsection}{\normalfont\bfseries\color{accent}\sffamily}
{\llap{\color{accent}\rule{1.2em}{1.2em}\hspace{0.5em}}}{0em}{}
\titleformat{\subsubsection}{\normalfont\bfseries\itshape\color{secondary}\sffamily}
{\llap{\color{secondary}\rule{1em}{1em}\hspace{0.5em}}}{0em}{}

% --- 6. MÔI TRƯỜNG ---
\def\phanI{TỔNG QUAN}
\def\phanII{CHUẨN BỊ CỦA GIÁO VIÊN}
\def\phanIII{KẾ HOẠCH DẠY HỌC ĐỀ XUẤT}
\def\phanIV{HƯỚNG DẪN HỌC SINH}
\def\phanV{CÁC PHỤ LỤC}

\newtcolorbox{competitionsection}[1]{
    enhanced, breakable, colback=white, colframe=primary, boxrule=0.8pt,
    fonttitle=\bfseries\sffamily\Huge, coltitle=white, colbacktitle=primary,
    attach boxed title to top left={xshift=1cm, yshift=-1mm-\tcboxedtitleheight/2},
    boxed title style={sharp corners, boxrule=0pt},
    title={PHẦN #1: \csname phan#1\endcsname},
    before upper={\addcontentsline{toc}{section}{Phần #1: \csname phan#1\endcsname}}
}

\newcolumntype{C}{>{\centering\arraybackslash}X}
\newcolumntype{L}{>{\raggedright\arraybackslash}X}
\newcommand{\header}[1]{\cellcolor{TableHeader}\sffamily\bfseries#1}

\newcommand{\teachingplan}[4]{
\noindent\begin{minipage}{0.12\textwidth}\centering\bfseries\sffamily\color{primary}#1\end{minipage}%
\begin{minipage}{0.23\textwidth}\sffamily\textbf{#2}\end{minipage}%
\begin{minipage}{0.35\textwidth}\textit{#3}\end{minipage}%
\begin{minipage}{0.25\textwidth}\small #4\end{minipage}
\vspace{1ex}\par\noindent\color{primary!30}\rule{\textwidth}{0.4pt}\vspace{1ex}}

\newtcolorbox{casestudy}{enhanced, breakable, colback=primary!5, colframe=primary!80, boxrule=1pt, arc=2mm,
    title=\sffamily\bfseries\color{accent}TÌNH HUỐNG THỰC TIỄN: STARTUP VẬN TẢI SMARTLOGI, fonttitle=\bfseries\sffamily}
\newtcolorbox{keyformula}{enhanced, breakable, colback=accent!5, colframe=accent!80, boxrule=1pt, arc=2mm,
    title=\sffamily\bfseries\color{primary}KIẾN THỨC TRỌNG TÂM, fonttitle=\bfseries\sffamily}
\newtcolorbox{task}[1]{enhanced, breakable, colback=warning!10, colframe=warning!80, boxrule=1pt, arc=2mm,
    title=\sffamily\bfseries\color{secondary}NHIỆM VỤ: #1, fonttitle=\bfseries\sffamily}
\newtcolorbox{phieuhoc_tap_box}[1]{enhanced, breakable, colback=white, colframe=success, boxrule=1pt, arc=2mm,
    title=\sffamily\bfseries #1, fonttitle=\bfseries\sffamily, coltitle=white}

\setlist[itemize,1]{label=\textbullet, leftmargin=*, nosep}
\setlist[itemize,2]{label=--, leftmargin=*, nosep}
\setlist[description]{font=\bfseries\sffamily\color{primary}, nosep}

\begin{document}

% --- TRANG BÌA ---
\begin{titlepage}
\begin{tikzpicture}[remember picture, overlay]
\node[rectangle, fill=primary, anchor=north, minimum width=\paperwidth, minimum height=2cm] (header) at (current page.north){};
\node[anchor=south, yshift=2mm, text=white] at (header.south) {\Huge\sffamily\bfseries GIÁO ÁN DỰ THI STEM/STEAM CẤP THÀNH PHỐ};
\end{tikzpicture}

\vspace{2cm}
\begin{center}
{\LARGE\bfseries\color{primary}\sffamily SỞ GIÁO DỤC VÀ ĐÀO TẠO THÀNH PHỐ HỒ CHÍ MINH}\\
\vspace{0.5cm}
{\Large\bfseries\sffamily HỘI THI THIẾT KẾ CHỦ ĐỀ DẠY HỌC STEM/STEAM}\\
\vspace{0.3cm}
{\large\sffamily Năm học 2025 - 2026}

\vspace{2cm}

\begin{tcolorbox}[width=0.9\textwidth, colback=gray!5, colframe=primary, arc=2mm, boxrule=1pt]
    \centering
    {\LARGE\bfseries\sffamily CHỦ ĐỀ:} \\
    {\Huge\bfseries\color{accent}\sffamily LOGISTICS THÔNG MINH}\\
    \vspace{0.25cm}
    {\Large\sffamily TỐI ƯU HÓA VẬN TẢI BẰNG HỆ BẤT PHƯƠNG TRÌNH}\\
    \vspace{0.25cm}
    {\large\sffamily Ứng dụng Quy hoạch tuyến tính và Công nghệ số}
\end{tcolorbox}

\vspace{1.5cm}
{\Large\bfseries\sffamily MÔN: TOÁN - KHỐI: 10}

\vfill
\large
\begin{tabular}{rl}
\sffamily\bfseries Giáo viên: & Nguyễn Văn Sang\\
\sffamily\bfseries Đơn vị: & Trường THPT Nguyễn Hữu Cảnh\\
\sffamily\bfseries Số điện thoại: & 0389\,821\,115\\
\end{tabular}

\vspace{1cm}
{\large\sffamily Thành phố Hồ Chí Minh, tháng 11 năm 2025}
\end{center}
\end{titlepage}

\tableofcontents
\thispagestyle{empty}

%%%%% PHẦN I: TỔNG QUAN %%%%%
\begin{competitionsection}{I}
\vspace{0.5cm}
\subsection*{1.1. Thông tin cá nhân}
\begin{itemize}[leftmargin=*,noitemsep]
\item \textbf{Họ và tên:} Nguyễn Văn Sang
\item \textbf{Chuyên môn:} Toán học
\item \textbf{Thâm niên:} 15 năm giảng dạy THPT
\item \textbf{Đơn vị:} Trường THPT Nguyễn Hữu Cảnh, TP.HCM
\end{itemize}

\subsection*{1.2. Giới thiệu chủ đề}
\begin{itemize}[leftmargin=*,noitemsep]
\item \textbf{Tên chủ đề:} \textcolor{accent}{Logistics Thông Minh: Tối ưu hóa bài toán vận tải bằng Hệ bất phương trình bậc nhất hai ẩn}
\item \textbf{Bối cảnh thực tiễn:}
    \begin{itemize}
    \item \textbf{82\%} doanh nghiệp vừa và nhỏ gặp khó khăn trong tối ưu chi phí vận chuyển (Khảo sát VCCI 2024)
    \item Sự bùng nổ của Thương mại điện tử (Shopee, TikTok Shop) đòi hỏi giải pháp logistics hiệu quả
    \item Học sinh cần thấy được ``sức mạnh kinh tế'' của Toán học trong kinh doanh thực tế
    \end{itemize}
\item \textbf{Điểm mới và sáng tạo:}
    \begin{itemize}
    \item \textbf{Đa ngành:} Kết hợp Toán học - Kinh tế - Công nghệ thông tin tạo thành vòng tròn STEAM hoàn chỉnh
    \item \textbf{Sản phẩm thực:} Công cụ web \texttt{SmartLogi.html} hoàn thiện có thể sử dụng ngay, không cần server
    \item \textbf{Tư duy khởi nghiệp:} Học sinh đóng vai CEO startup để giải quyết bài toán lỗ/lãi thực tế
    \item \textbf{Trực quan hóa:} Biến các con số khô khan thành đồ thị miền nghiệm sinh động
    \end{itemize}
\end{itemize}

\subsection*{1.3. Căn cứ thiết kế}
\begin{minipage}{0.48\textwidth}
\textbf{Căn cứ pháp lý:}
\begin{itemize}[leftmargin=*,noitemsep]
\item Chương trình GDPT 2018 môn Toán lớp 10
\item Công văn 5555/BGDĐT-GDTrH về dạy học STEM
\item Kế hoạch số 4500/KH-SGDĐT của Sở GD\&ĐT TP.HCM
\end{itemize}
\end{minipage}
\hfill
\begin{minipage}{0.48\textwidth}
\textbf{Cơ sở lý luận:}
\begin{itemize}[leftmargin=*,noitemsep]
\item Lý thuyết học tập qua dự án (Project-Based Learning)
\item Mô hình 5E (Engage, Explore, Explain, Elaborate, Evaluate)
\item Khung năng lực khởi nghiệp của EU (EntreComp)
\end{itemize}
\end{minipage}

\subsection*{1.4. Mục tiêu giáo dục}
\begin{center}
\begin{tabularx}{\linewidth}{@{} L L L @{}}
\toprule
\header{Kiến thức} & \header{Kỹ năng} & \header{Thái độ} \\
\midrule
\parbox[t]{\dimexpr1.\linewidth-2\tabcolsep}{
    \begin{itemize}[leftmargin=*,noitemsep,topsep=1ex]
        \item Xác định miền nghiệm của Hệ BPT bậc nhất 2 ẩn
        \item Hiểu bản chất bài toán tối ưu (Linear Programming)
        \item Tìm cực trị của hàm $F(x,y)$ trên miền đa giác
    \end{itemize}
} &
\parbox[t]{\dimexpr1.\linewidth-2\tabcolsep}{
    \begin{itemize}[leftmargin=*,noitemsep,topsep=1ex]
        \item Mô hình hóa bài toán thực tế
        \item Sử dụng công nghệ số (Web App)
        \item Đọc hiểu đồ thị, phân tích kịch bản
        \item Làm việc nhóm, thuyết trình
    \end{itemize}
} &
\parbox[t]{\dimexpr1.\linewidth-2\tabcolsep}{
    \begin{itemize}[leftmargin=*,noitemsep,topsep=1ex]
        \item Tư duy tối ưu hóa (chi phí thấp nhất)
        \item Tinh thần khởi nghiệp
        \item Trách nhiệm với môi trường (giảm khí thải)
    \end{itemize}
} \\
\bottomrule
\end{tabularx}
\end{center}

\subsection*{1.5. Cấu trúc chương trình}
\begin{center}
\begin{tikzpicture}[node distance=1.5cm, every node/.style={font=\sffamily}]
\node (start) [rectangle, rounded corners, fill=primary!20, text width=4.5cm, align=center, minimum height=1.2cm] {\textbf{Khám phá thực tiễn}\\ CEO Startup cần tối ưu chi phí};
\node (mid1) [rectangle, rounded corners, fill=accent!20, text width=4.5cm, align=center, below of=start, yshift=-1.2cm, minimum height=1.2cm] {\textbf{Xây dựng mô hình Toán}\\ Hệ BPT, Miền nghiệm};
\node (mid2) [rectangle, rounded corners, fill=primary!20, text width=4.5cm, align=center, below of=mid1, yshift=-1.2cm, minimum height=1.2cm] {\textbf{Phát triển công cụ số}\\ Ứng dụng SmartLogi.html};
\node (end) [rectangle, rounded corners, fill=accent!20, text width=4.5cm, align=center, below of=mid2, yshift=-1.2cm, minimum height=1.2cm] {\textbf{Ứng dụng thực tế}\\ Lập kế hoạch vận tải tối ưu};
\draw[->, thick, color=primary] (start) -- (mid1) node[midway, right, text width=3cm, align=center] {\footnotesize \textbf{M} (Toán học)};
\draw[->, thick, color=primary] (mid1) -- (mid2) node[midway, right, text width=3cm, align=center] {\footnotesize \textbf{T} (Công nghệ)};
\draw[->, thick, color=primary] (mid2) -- (end) node[midway, right, text width=3cm, align=center] {\footnotesize \textbf{A} (Kinh tế)};
\end{tikzpicture}
\end{center}
\end{competitionsection}
\newpage

%%%%% PHẦN II: CHUẨN BỊ %%%%%
\begin{competitionsection}{II}
\vspace{0.5cm}
\subsection*{2.1. Thiết bị và công nghệ}
\begin{itemize}[leftmargin=*,noitemsep]
\item \textbf{Phần cứng:}
    \begin{itemize}
    \item Máy tính có kết nối Internet (1 máy/2 HS) hoặc điện thoại thông minh
    \item Máy chiếu tương tác thông minh
    \end{itemize}
\item \textbf{Phần mềm và nền tảng:}
    \begin{itemize}
    \item \textbf{Học liệu số chính:} Ứng dụng web \texttt{SmartLogi.html} (đính kèm) hoặc truy cập \href{https://stem-1h8.pages.dev/SmartLogi.html}{stem-1h8.pages.dev/SmartLogi.html}
    \item \textit{Yêu cầu:} Trình duyệt web (Chrome, Firefox, Edge...)
    \item Google Sheets template (dự phòng)
    \item Padlet để chia sẻ kết quả nhóm
    \end{itemize}
\item \textbf{Tài liệu:}
    \begin{itemize}
    \item Phiếu học tập số 1: Khám phá bài toán tối ưu
    \item Phiếu học tập số 2: Hướng dẫn sử dụng 3 Tab của SmartLogi
    \item Phiếu khảo sát trước/sau học (đánh giá nhận thức học sinh)
    \item Bảng tham khảo giá thuê xe thực tế
    \item Rubric đánh giá sản phẩm
    \item Phiếu đánh giá đồng cấp (Peer Assessment)
    \end{itemize}
\end{itemize}

\subsection*{2.2. Kịch bản giảng dạy}
\begin{center}
\begin{tabularx}{\linewidth}{@{} L L L L L @{}}
\toprule
\header{Thời lượng} & \header{Hoạt động} & \header{Phương pháp} & \header{Công cụ} & \header{Đánh giá} \\
\midrule
10' & Khởi động: Game CEO & Tình huống thực tiễn & Clip + Web App & Quan sát hứng thú \\
15' & Khám phá Hệ BPT & Khám phá có hướng dẫn & PHT số 1 & Phiếu trả lời \\
\rowcolor{primary!10}
10' & Chuyển tiếp công nghệ & Đàm thoại & Demo App (3 tabs) & Câu hỏi phản hồi \\
\rowcolor{primary!10}
25' & Thực hành (Tab 1, 2) & Học qua dự án & Web App SmartLogi & Đánh giá quá trình \\
\rowcolor{primary!10}
15' & Phân tích kết quả & Hợp tác nhóm & Biểu đồ (Tab 2) & Peer assessment \\
10' & Trình bày sản phẩm & Thuyết trình & Google Slides & Rubric đánh giá \\
5' & Tổng kết mở rộng & Suy ngẫm & Tab 3: Học tập & Nhật ký học tập \\
\bottomrule
\end{tabularx}
\end{center}

\subsection*{2.3. Đánh giá đa chiều}
\begin{minipage}[t]{0.48\textwidth}
\textbf{Đánh giá quá trình:}
\begin{itemize}[leftmargin=*,noitemsep]
\item \textbf{Phiếu quan sát nhóm:} Mức độ hợp tác, giải quyết vấn đề
\item \textbf{Nhật ký học tập:} Suy nghĩ, thắc mắc của học sinh
\item \textbf{Đánh giá đồng cấp:} Học sinh đánh giá lẫn nhau
\end{itemize}
\end{minipage}
\hfill
\begin{minipage}[t]{0.48\textwidth}
\textbf{Đánh giá sản phẩm:}
\begin{itemize}[leftmargin=*,noitemsep]
\item \textbf{Rubric 4 tiêu chí:} (Xem Phụ lục 3)
    \begin{itemize}
    \item Tính chính xác mô hình toán (4 điểm)
    \item Khả năng sử dụng công nghệ (4 điểm)
    \item Tính khả thi phương án (4 điểm)
    \item Khả năng trình bày (4 điểm)
    \end{itemize}
\end{itemize}
\end{minipage}

\subsection*{2.4. Kế hoạch dự phòng}
\begin{itemize}[leftmargin=*,noitemsep]
\item \textbf{Tình huống 1:} Mất điện hoặc mất kết nối Internet
    \begin{itemize}
    \item Học liệu số \texttt{SmartLogi.html} \textbf{không cần Internet}, chỉ cần mở bằng trình duyệt.
    \item Nếu mất điện: Chuyển sang sử dụng Google Sheets offline hoặc phiếu tính toán bằng tay.
    \end{itemize}
\item \textbf{Tình huống 2:} Học sinh không có thiết bị cá nhân
    \begin{itemize}
    \item Sắp xếp nhóm 3-4 học sinh/1 máy tính của trường
    \item Sử dụng phương pháp ``station rotation'' để luân phiên sử dụng thiết bị
    \end{itemize}
\item \textbf{Tình huống 3:} Một số học sinh tiếp thu nhanh hơn
    \begin{itemize}
    \item Giao nhiệm vụ mở rộng: Khám phá ``Tab Học tập'' và ``Phân tích độ nhạy''.
    \item Giao vai trò hỗ trợ các nhóm khác
    \end{itemize}
\end{itemize}
\end{competitionsection}

%%%%% PHẦN III: KẾ HOẠCH DẠY HỌC %%%%%
\begin{competitionsection}{III}
\vspace{0.5cm}
\subsection*{3.1. Tiến trình hai tiết học (Tóm tắt)}
\teachingplan{10'}{Khởi động}{Video tình huống: CEO startup cần tối ưu chi phí vận chuyển}{Công nghệ: Clip tương tác}
\teachingplan{15'}{Khám phá toán học}{Xây dựng Hệ BPT từ bài toán thực tế}{PP: Khám phá có hướng dẫn}
\teachingplan{10'}{Chuyển tiếp}{Demo ứng dụng web SmartLogi (3 tabs)}{Công nghệ: Màn hình tương tác}
\teachingplan{25'}{Thực hành}{Sử dụng Tab 1 (Tính toán) và Tab 2 (Kịch bản)}{PP: Học qua dự án}
\teachingplan{15'}{Phân tích}{So sánh kịch bản, phân tích độ nhạy}{Công cụ: Biểu đồ cột, đường}
\teachingplan{10'}{Báo cáo}{Thuyết trình nhóm: Phương án tối ưu}{Công nghệ: Google Slides chia sẻ}
\teachingplan{5'}{Tổng kết}{Hệ thống hóa kiến thức (Tab 3)}{PP: Sơ đồ tư duy}

\subsection*{3.2. Kịch bản chi tiết}

\begin{casestudy}
Bạn là CEO của startup SmartLogi - một công ty logistics mới thành lập. Công ty nhận được đơn hàng đầu tiên: vận chuyển tối thiểu \textbf{9 tấn} hàng hóa với tổng thể tích tối thiểu \textbf{36 m$^3$}. Bạn có thể thuê 2 loại xe:
\begin{itemize}[leftmargin=*,noitemsep]
\item \textbf{Xe loại A:} Tải trọng 3 tấn, thể tích 3 m$^3$, giá thuê 4 triệu/xe
\item \textbf{Xe loại B:} Tải trọng 1 tấn, thể tích 6 m$^3$, giá thuê 3 triệu/xe
\end{itemize}
\textbf{Thách thức:} Với ngân sách eo hẹp của startup, làm thế nào để chọn phương án vận chuyển với chi phí thấp nhất mà vẫn đảm bảo yêu cầu đơn hàng?
\end{casestudy}

\subsubsection*{Hoạt động 1: Khởi động - Kết nối thực tiễn (10 phút)}
\begin{itemize}[leftmargin=*,noitemsep]
\item \textbf{Bước 1:} Giáo viên chiếu video giới thiệu về ngành Logistics và thương mại điện tử tại Việt Nam
\item \textbf{Bước 2:} Đặt câu hỏi kích thích tư duy:
    \begin{itemize}
    \item \textit{``Nếu em là CEO một startup logistics, em sẽ làm gì để tối ưu chi phí?''}
    \item \textit{``Làm thế nào để biết nên thuê bao nhiêu xe mỗi loại?''}
    \end{itemize}
\item \textbf{Bước 3:} Học sinh thảo luận cặp đôi (2 phút) và chia sẻ ý tưởng ban đầu
\item \textbf{Dự kiến câu trả lời:}
    \begin{itemize}
    \item Thuê nhiều xe nhỏ để linh hoạt
    \item Thuê xe lớn để tiết kiệm chi phí
    \item Cần tính toán cụ thể để so sánh
    \end{itemize}
\item \textbf{Chuyển tiếp:} \textit{``Hôm nay chúng ta sẽ sử dụng Toán học và Công nghệ để giải quyết bài toán kinh doanh thực tế này!''}
\end{itemize}

\subsubsection*{Hoạt động 2: Khám phá Hệ bất phương trình (15 phút)}
\begin{task}{Nhóm}
\textbf{Nhiệm vụ:} Mô hình hóa bài toán logistics thành Hệ bất phương trình (Phiếu học tập số 1)

\textbf{Hướng dẫn:}
\begin{itemize}[leftmargin=*,noitemsep]
\item Xác định các ẩn số: $x$ (số xe A), $y$ (số xe B)
\item Viết các ràng buộc từ yêu cầu bài toán
\item Xác định hàm mục tiêu: Chi phí cần tối thiểu hóa
\end{itemize}
\end{task}

\begin{keyformula}
\textbf{Hệ bất phương trình bậc nhất hai ẩn:}
$$\begin{cases}
3x + y \ge 9 & \text{(Ràng buộc tải trọng)} \\
3x + 6y \ge 36 & \text{(Ràng buộc thể tích)} \\
0 \le x \le 10 & \text{(Giới hạn xe A)} \\
0 \le y \le 9 & \text{(Giới hạn xe B)}
\end{cases}$$

\textbf{Hàm mục tiêu:} $C = 4x + 3y \to \min$

\textbf{Ý nghĩa:}
\begin{itemize}[leftmargin=*,noitemsep]
\item \textbf{Miền nghiệm:} Tập hợp các điểm $(x, y)$ thỏa mãn tất cả các ràng buộc
\item \textbf{Điểm tối ưu:} Điểm trong miền nghiệm làm cho hàm $C$ đạt giá trị nhỏ nhất
\item \textbf{Định lý:} Cực trị của hàm tuyến tính trên miền đa giác lồi đạt được tại đỉnh
\end{itemize}
\end{keyformula}

\subsubsection*{Hoạt động 3: Chuyển tiếp sang công nghệ (10 phút)}
\begin{itemize}[leftmargin=*,noitemsep]
\item \textbf{Bước 1:} Giáo viên demo ứng dụng web \texttt{SmartLogi.html} (Phiếu học tập số 2).
\item \textbf{Bước 2:} Thao tác minh họa 3 tab chính:
    \begin{itemize}
    \item \textbf{Tab 1 - Tính toán:} Nhập các tham số xe A, xe B, yêu cầu tải trọng/thể tích. Cho thấy miền nghiệm (vùng tô màu) và điểm tối ưu (điểm đỏ).
    \item \textbf{Tab 2 - Phân tích kịch bản:} Thêm 2-3 kịch bản (VD: thay đổi giá xe, thay đổi yêu cầu). Cho thấy biểu đồ cột so sánh chi phí. Kéo xuống xem biểu đồ ``Phân tích độ nhạy''.
    \item \textbf{Tab 3 - Học tập:} Lướt qua nội dung công thức Toán, kiến thức Logistics và quy trình STEM.
    \end{itemize}
\item \textbf{Bước 3:} Đặt câu hỏi chuyển tiếp:
    \begin{itemize}
    \item \textit{``Tại sao điểm tối ưu lại nằm ở đỉnh của miền nghiệm?''}
    \item \textit{``Nếu giá xe thay đổi, điểm tối ưu có thay đổi không?''}
    \end{itemize}
\end{itemize}

\subsubsection*{Hoạt động 4: Thực hành với Web App (25 phút)}
\begin{task}{Nhóm}
\textbf{Nhiệm vụ:} Sử dụng Web App SmartLogi để tìm phương án tối ưu và phân tích các kịch bản.

\textbf{Yêu cầu cụ thể:}
\begin{itemize}[leftmargin=*,noitemsep]
\item \textbf{Tab 1:} Nhập kịch bản gốc theo Phiếu học tập số 1, ghi lại kết quả tối ưu.
\item \textbf{Tab 2:} Sử dụng nút ``Thêm kịch bản'' để thử nghiệm các thay đổi:
    \begin{itemize}
    \item Kịch bản 2: Tăng giá xe A lên 5 triệu
    \item Kịch bản 3: Tăng yêu cầu tải trọng lên 12 tấn
    \item Kịch bản 4: Nhóm tự đề xuất
    \end{itemize}
\item \textbf{Tab 2:} Quan sát biểu đồ ``Phân tích độ nhạy'' và nhận xét.
\item \textbf{Chụp màn hình} kết quả để chuẩn bị báo cáo.
\end{itemize}
\end{task}

\begin{itemize}[leftmargin=*,noitemsep]
\item \textbf{Hỗ trợ giáo viên:}
    \begin{itemize}
    \item Tuần tra các nhóm, hỗ trợ kỹ thuật khi cần
    \item Gợi ý cách đọc đồ thị miền nghiệm
    \item Hướng dẫn cách đọc biểu đồ cột (so sánh) và biểu đồ đường (độ nhạy).
    \end{itemize}
\end{itemize}

\subsubsection*{Hoạt động 5: Phân tích và trình bày (15 phút)}
\begin{itemize}[leftmargin=*,noitemsep]
\item \textbf{Bước 1:} 2 nhóm được chọn ngẫu nhiên trình bày phương án của mình (3 phút/nhóm).
\item \textbf{Bước 2:} Các nhóm khác đặt câu hỏi phản biện
\item \textbf{Bước 3:} Giáo viên đặt câu hỏi mở rộng:
    \begin{itemize}
    \item \textit{``Dựa trên biểu đồ 'Phân tích độ nhạy', yếu tố nào (giá xe hay yêu cầu tải trọng) có ảnh hưởng mạnh hơn đến chi phí?''}
    \item \textit{``Nếu khách hàng đột ngột tăng yêu cầu thể tích lên 50 m$^3$, phương án của nhóm có còn khả thi không?''}
    \end{itemize}
\item \textbf{Bước 4:} Cả lớp thảo luận và thống nhất phương án tối ưu
\end{itemize}

\subsubsection*{Hoạt động 6: Tổng kết và mở rộng (5 phút)}
\begin{itemize}[leftmargin=*,noitemsep]
\item \textbf{Tổng kết kiến thức:} Giáo viên hệ thống hóa (sử dụng \textbf{Tab 3 - Học tập}):
    \begin{itemize}
    \item Hệ BPT bậc nhất 2 ẩn giúp mô hình hóa bài toán kinh doanh thực tế (M).
    \item Công nghệ (T) giúp trực quan hóa miền nghiệm và tự động tìm điểm tối ưu.
    \item Tư duy tối ưu hóa là kỹ năng quan trọng trong kinh doanh (A - Arts/Kinh tế).
    \end{itemize}
\item \textbf{Bài tập về nhà:}
    \begin{itemize}
    \item Sử dụng công cụ SmartLogi để giải quyết một bài toán logistics khác (VD: cho công ty gia đình hoặc cửa hàng online).
    \item Trình bày trong 1 trang giấy hoặc slide, bao gồm:
        \begin{itemize}
        \item Bài toán cụ thể (yêu cầu, ràng buộc)
        \item Phương án tối ưu được lựa chọn (ảnh chụp từ Tab 1, 2)
        \item Lý do lựa chọn và tính khả thi
        \end{itemize}
    \end{itemize}
\end{itemize}
\end{competitionsection}
\newpage

%%%%% PHẦN IV: HƯỚNG DẪN HỌC SINH %%%%%
\begin{competitionsection}{IV}
\vspace{0.5cm}
\subsection*{4.1. Hướng dẫn chi tiết cho học sinh}
\subsubsection*{Giai đoạn 1: Mô hình hóa bài toán (Phiếu học tập số 1)}
\begin{itemize}[leftmargin=*,noitemsep]
\item \textbf{Bước 1:} Đọc kỹ tình huống CEO startup SmartLogi
\item \textbf{Bước 2:} Xác định các ẩn số:
    \begin{itemize}
    \item $x$ = số xe loại A cần thuê
    \item $y$ = số xe loại B cần thuê
    \end{itemize}
\item \textbf{Bước 3:} Viết các ràng buộc từ yêu cầu bài toán:
    \begin{itemize}
    \item Tải trọng: $3x + 1y \ge 9$ (cần ít nhất 9 tấn)
    \item Thể tích: $3x + 6y \ge 36$ (cần ít nhất 36 m$^3$)
    \item Giới hạn: $0 \le x \le 10$, $0 \le y \le 9$
    \end{itemize}
\item \textbf{Bước 4:} Viết hàm mục tiêu: $C = 4x + 3y \to \min$
\item \textbf{Bước 5:} Dự đoán trực giác và kiểm tra tính khả thi
\end{itemize}

\subsubsection*{Giai đoạn 2: Sử dụng công cụ SmartLogi (Phiếu học tập số 2)}
\begin{itemize}[leftmargin=*,noitemsep]
\item \textbf{Bước 1: Mở ứng dụng}
    \begin{itemize}
    \item Mở file \texttt{SmartLogi.html} bằng trình duyệt web.
    \end{itemize}
\item \textbf{Bước 2: Tab 1 - Tính toán}
    \begin{itemize}
    \item Nhập các tham số của bài toán (xe A, xe B, yêu cầu tải trọng/thể tích).
    \item Quan sát miền nghiệm (vùng tô màu xanh) trên đồ thị.
    \item Ghi lại điểm tối ưu (x, y) và chi phí tối thiểu.
    \end{itemize}
\item \textbf{Bước 3: Tab 2 - Phân tích kịch bản}
    \begin{itemize}
    \item Nhấn nút ``Thêm kịch bản'' để tạo phương án mới.
    \item Thay đổi các tham số (VD: Tăng giá xe, thay đổi yêu cầu).
    \item Quan sát biểu đồ cột để so sánh các kịch bản.
    \item Kéo xuống, quan sát biểu đồ đường ``Phân tích độ nhạy''.
    \end{itemize}
\item \textbf{Bước 4: Tab 3 - Học tập}
    \begin{itemize}
    \item Đọc lại công thức Hệ BPT và điều kiện tối ưu.
    \item Đọc phần ``Kiến thức Logistics'' để hiểu thêm về ngành.
    \item Đọc phần ``Quy trình STEM'' để hiểu rõ dự án.
    \end{itemize}
\end{itemize}

\subsubsection*{Giai đoạn 3: Phân tích và trình bày}
\begin{itemize}[leftmargin=*,noitemsep]
\item \textbf{Bước 1:} So sánh các phương án (từ Tab 2) dựa trên tiêu chí:
    \begin{itemize}
    \item Chi phí thấp nhất
    \item Tính linh hoạt (dễ điều chỉnh khi yêu cầu thay đổi)
    \item Rủi ro (độ nhạy với biến động giá)
    \end{itemize}
\item \textbf{Bước 2:} Chọn 1 phương án tối ưu và lý giải lựa chọn
\item \textbf{Bước 3:} Chuẩn bị phần trình bày 3 phút với cấu trúc:
    \begin{itemize}
    \item Giới thiệu bài toán và vai trò CEO startup
    \item Trình bày các phương án đã thử nghiệm (chụp ảnh biểu đồ từ Tab 2)
    \item Phân tích và lựa chọn phương án tối ưu
    \item Kết luận (có nhắc đến độ nhạy và rủi ro)
    \end{itemize}
\end{itemize}

\subsection*{4.2. Gợi ý mở rộng}
\begin{itemize}[leftmargin=*,noitemsep]
\item \textbf{Tìm hiểu sâu về Logistics (từ Tab 3):}
    \begin{itemize}
    \item Nghiên cứu các thuật toán tối ưu hóa phức tạp hơn (Simplex, Branch and Bound)
    \item Tìm hiểu về chuỗi cung ứng (Supply Chain) và quản lý kho
    \item Khám phá các phần mềm logistics chuyên nghiệp
    \end{itemize}
\item \textbf{Phát triển công cụ (từ Phụ lục 5):}
    \begin{itemize}
    \item Thử ``View Source'' (Xem nguồn trang) file HTML để đọc code JavaScript.
    \item Đề xuất cải tiến giao diện hoặc thuật toán tính toán
    \item Thêm tính năng mới (VD: tính toán thời gian, tuyến đường)
    \end{itemize}
\item \textbf{Ứng dụng thực tế:}
    \begin{itemize}
    \item Giúp gia đình hoặc cửa hàng nhỏ tối ưu vận chuyển hàng hóa
    \item Tham gia cuộc thi khởi nghiệp học sinh do trường tổ chức
    \item Viết bài chia sẻ kinh nghiệm trên blog cá nhân hoặc mạng xã hội
    \end{itemize}
\end{itemize}
\end{competitionsection}
\newpage

%%%%% PHẦN V: PHỤ LỤC %%%%%
\begin{competitionsection}{V}
\vspace{0.5cm}
\subsection*{Phụ lục 1: Phiếu học tập số 1 - Khám phá bài toán tối ưu}
\begin{phieuhoc_tap_box}{Phụ lục 1: Phiếu học tập số 1 - Mô hình hóa bài toán Logistics}
\textbf{Tình huống:} Bạn là CEO startup SmartLogi, cần vận chuyển tối thiểu 9 tấn hàng và 36 m$^3$ hàng.

\vspace{0.5cm}
\textbf{Bài 1: Xác định ẩn số}
\begin{itemize}[leftmargin=*,noitemsep]
\item Gọi $x$ = số xe loại A cần thuê
\item Gọi $y$ = số xe loại B cần thuê
\item Điều kiện: $x, y \ge 0$ và $x, y \in \mathbb{N}$ (số nguyên)
\end{itemize}

\textbf{Bài 2: Viết các ràng buộc (điền vào chỗ trống)}
\begin{itemize}[leftmargin=*,noitemsep]
\item Tải trọng (Xe A chở 3 tấn, Xe B chở 1 tấn, cần $\ge$ 9 tấn): $3x + \underline{\hspace{1cm}}y \ge \underline{\hspace{1cm}}$
\item Thể tích (Xe A chở 3 m$^3$, Xe B chở 6 m$^3$, cần $\ge$ 36 m$^3$): $\underline{\hspace{1cm}}x + 6y \ge \underline{\hspace{1cm}}$
\item Giới hạn xe A: $0 \le x \le \underline{\hspace{1cm}}$
\item Giới hạn xe B: $0 \le y \le \underline{\hspace{1cm}}$
\end{itemize}

\textbf{Bài 3: Xác định hàm mục tiêu}
\begin{itemize}[leftmargin=*,noitemsep]
\item Chi phí thuê xe A: 4 triệu/xe
\item Chi phí thuê xe B: 3 triệu/xe
\item Hàm chi phí: $C = \underline{\hspace{1cm}}x + \underline{\hspace{1cm}}y \to \min$
\end{itemize}

\textbf{Bài 4: Dự đoán theo cảm tính}
\begin{itemize}[leftmargin=*,noitemsep]
\item Phương án dự đoán của em: \underline{\hspace{2cm}} xe A, \underline{\hspace{2cm}} xe B
\item Chi phí dự đoán: \underline{\hspace{3cm}} triệu VNĐ
\item Kiểm tra: Tải trọng = \underline{\hspace{2cm}} tấn $\ge$ 9 tấn? \quad (\text{Đạt} / \text{Không đạt})
\item Kiểm tra: Thể tích = \underline{\hspace{2cm}} m$^3$ $\ge$ 36 m$^3$? \quad (\text{Đạt} / \text{Không đạt})
\end{itemize}
\end{phieuhoc_tap_box}

\newpage
\subsection*{Phụ lục 2: Phiếu học tập số 2 - Hướng dẫn sử dụng SmartLogi}
\begin{phieuhoc_tap_box}{Phụ lục 2: Phiếu học tập số 2 - Nhiệm vụ với Web App SmartLogi.html}
\textbf{\color{accent}I. Giới thiệu ứng dụng}
\begin{itemize}[leftmargin=*,noitemsep]
\item \textbf{Tên ứng dụng:} SmartLogi - Tối ưu hóa Vận tải
\item \textbf{Địa chỉ truy cập:} File \texttt{SmartLogi.html} (mở bằng trình duyệt) hoặc \href{https://stem-1h8.pages.dev/SmartLogi.html}{stem-1h8.pages.dev/SmartLogi.html}
\item \textbf{Cấu trúc 3 Tab:}
    \begin{itemize}
    \item \textbf{Tab 1 - TÍNH TOÁN:} Nhập tham số, xem miền nghiệm, tìm điểm tối ưu.
    \item \textbf{Tab 2 - PHÂN TÍCH KỊCH BẢN:} Thêm/xóa kịch bản, so sánh biểu đồ cột, phân tích độ nhạy.
    \item \textbf{Tab 3 - HỌC TẬP:} Công thức Toán, kiến thức Logistics, quy trình STEM.
    \end{itemize}
\end{itemize}

\textbf{\color{accent}II. Hướng dẫn thực hiện nhiệm vụ}
\begin{itemize}[leftmargin=*,noitemsep]
\item \textbf{Bước 1: Mở Tab ``TÍNH TOÁN''}
    \begin{itemize}
    \item Nhập các tham số của bài toán gốc:
    \item \textit{Xe A:} Tải trọng = 3 tấn, Thể tích = 3 m$^3$, Giá = 4 triệu
    \item \textit{Xe B:} Tải trọng = 1 tấn, Thể tích = 6 m$^3$, Giá = 3 triệu
    \item \textit{Yêu cầu:} Tổng tấn = 9, Tổng khối = 36
    \item \textit{Giới hạn:} Max xe A = 10, Max xe B = 9
    \item \textit{Ghi lại kết quả tối ưu:} \underline{\hspace{2cm}} xe A, \underline{\hspace{2cm}} xe B
    \item \textit{Chi phí tối ưu:} \underline{\hspace{4cm}} triệu
    \end{itemize}
\item \textbf{Bước 2: Mở Tab ``PHÂN TÍCH KỊCH BẢN''}
    \begin{itemize}
    \item Nhấn ``+ Thêm kịch bản'' để tạo \textbf{Kịch bản 2} của nhóm em
    \item \textit{Thay đổi:} (VD: Giá xe A tăng lên 5 triệu) \underline{\hspace{6cm}}
    \item \textit{Kết quả mới:} \underline{\hspace{2cm}} xe A, \underline{\hspace{2cm}} xe B. Chi phí: \underline{\hspace{3cm}}
    \item Quan sát biểu đồ cột. Kịch bản nào tốt hơn? \underline{\hspace{4cm}}
    \end{itemize}
\item \textbf{Bước 3: Phân tích Độ nhạy (vẫn ở Tab 2)}
    \begin{itemize}
    \item Kéo xuống xem biểu đồ đường ``Phân tích độ nhạy''
    \item \textit{Nhận xét:} Yếu tố nào làm chi phí thay đổi nhiều nhất? \underline{\hspace{8cm}}
    \end{itemize}
\item \textbf{Bước 4: Mở Tab ``HỌC TẬP''}
    \begin{itemize}
    \item Đọc phần ``Kiến thức Logistics''. Kể tên 2 thuật ngữ mới:
    \item 1: \underline{\hspace{4cm}} | 2: \underline{\hspace{4cm}}
    \end{itemize}
\end{itemize}
\end{phieuhoc_tap_box}

\newpage
\subsection*{Phụ lục 3: Rubric đánh giá sản phẩm}
\begin{tabularx}{\linewidth}{@{} >{\raggedright\arraybackslash}p{3.2cm} >{\raggedright\arraybackslash}X >{\raggedright\arraybackslash}X >{\raggedright\arraybackslash}X @{}}
    \toprule
    \rowcolor{TableHeader}
    \sffamily\bfseries Tiêu chí & \sffamily\bfseries Xuất sắc (4đ) & \sffamily\bfseries Khá (3đ) & \sffamily\bfseries Cần cải thiện (1-2đ) \\
    \midrule
    
    \textbf{Mô hình hóa Toán học} 
    & 
    \begin{itemize}[leftmargin=*,noitemsep,topsep=1ex]
        \item Hệ BPT chính xác
        \item Xác định đúng hàm mục tiêu
        \item Giải thích rõ ràng
    \end{itemize}
    & 
    \begin{itemize}[leftmargin=*,noitemsep,topsep=1ex]
        \item Hệ BPT cơ bản đúng
        \item Có sai sót nhỏ
    \end{itemize}
    & 
    \begin{itemize}[leftmargin=*,noitemsep,topsep=1ex]
        \item Hệ BPT sai
        \item Không xác định được mục tiêu
    \end{itemize} \\
    \addlinespace
    
    \textbf{Sử dụng công nghệ (SmartLogi App)} 
    & 
    \begin{itemize}[leftmargin=*,noitemsep,topsep=1ex]
        \item Thành thạo cả 3 tabs
        \item Phân tích độ nhạy đúng
    \end{itemize}
    & 
    \begin{itemize}[leftmargin=*,noitemsep,topsep=1ex]
        \item Sử dụng tốt Tab 1, 2
        \item Còn lúng túng Tab 3
    \end{itemize}
    & 
    \begin{itemize}[leftmargin=*,noitemsep,topsep=1ex]
        \item Chỉ dùng được Tab 1
        \item Không phân tích được
    \end{itemize} \\
    \addlinespace
    
    \textbf{Tính khả thi phương án} 
    & 
    \begin{itemize}[leftmargin=*,noitemsep,topsep=1ex]
        \item Đề xuất 2+ kịch bản khả thi
        \item Phân tích ưu nhược chi tiết
    \end{itemize}
    & 
    \begin{itemize}[leftmargin=*,noitemsep,topsep=1ex]
        \item Đề xuất 1-2 kịch bản
        \item Phân tích cơ bản
    \end{itemize}
    & 
    \begin{itemize}[leftmargin=*,noitemsep,topsep=1ex]
        \item Phương án không khả thi
        \item Thiếu phân tích
    \end{itemize} \\
    \addlinespace
    
    \textbf{Trình bày và bảo vệ} 
    & 
    \begin{itemize}[leftmargin=*,noitemsep,topsep=1ex]
        \item Mạch lạc, logic
        \item Trực quan hiệu quả
        \item Trả lời trọn vẹn
    \end{itemize}
    & 
    \begin{itemize}[leftmargin=*,noitemsep,topsep=1ex]
        \item Trình bày rõ ràng
        \item Có hỗ trợ trực quan
    \end{itemize}
    & 
    \begin{itemize}[leftmargin=*,noitemsep,topsep=1ex]
        \item Trình bày rời rạc
        \item Không trả lời được
    \end{itemize} \\
    \bottomrule
\end{tabularx}
\vspace{1em}
\noindent\textbf{Tổng điểm:} \underline{\hspace{2cm}}/16 điểm\\
\textbf{Nhận xét của giáo viên:}\\
\vspace{2cm}

\subsection*{Phụ lục 4: Bảng tham khảo giá thuê xe}
\begin{center}
\begin{tabularx}{0.95\linewidth}{@{} l C C C C @{}}
\toprule
\header{Loại xe} & \header{Tải trọng (tấn)} & \header{Thể tích (m$^3$)} & \header{Giá thuê (triệu/chuyến)} & \header{Ghi chú} \\
\midrule
Xe tải nhỏ (1 tấn) & 1.0 & 6 & 2.5 - 3.5 & Linh hoạt nội thành \\
Xe tải trung (3 tấn) & 3.0 & 12 & 4.0 - 5.0 & Phổ biến nhất \\
Xe tải lớn (5 tấn) & 5.0 & 20 & 6.0 - 7.5 & Đường dài \\
Xe van (chở đồ cồng kềnh) & 0.8 & 8 & 2.0 - 3.0 & Đồ nội thất \\
\bottomrule
\multicolumn{5}{@{}l}{\footnotesize * Số liệu tham khảo tháng 11/2025, giá có thể thay đổi theo khu vực và thời điểm.}
\end{tabularx}
\end{center}

\subsection*{Phụ lục 5: Mã nguồn ứng dụng (tham khảo)}
\textbf{Giới thiệu:} Ứng dụng SmartLogi được phát triển bằng \textbf{HTML, CSS và JavaScript (vanilla JS)}.
\begin{itemize}[leftmargin=*,noitemsep]
\item \textbf{Cấu trúc 3 Tabs:} (1) Tính toán \& Miền nghiệm, (2) Phân tích kịch bản, (3) Học tập.
\item \textbf{Thư viện:} \texttt{Chart.js} (vẽ biểu đồ), \texttt{SweetAlert2} (thông báo).
\item \textbf{Đặc điểm:} Chạy trực tiếp từ file \texttt{.html}, không cần server.
\end{itemize}

\textbf{Thuật toán tối ưu (JavaScript):}
\begin{verbatim}
// Brute-force Integer Linear Programming
function findOptimal(constraints, objective) {
    let minCost = Infinity;
    let optimalPoint = null;
    
    for (let x = 0; x <= maxA; x++) {
        for (let y = 0; y <= maxB; y++) {
            // Kiểm tra ràng buộc
            if (loadA*x + loadB*y >= minLoad &&
                volA*x + volB*y >= minVol) {
                // Tính chi phí
                const cost = costA*x + costB*y;
                if (cost < minCost) {
                    minCost = cost;
                    optimalPoint = {x, y, cost};
                }
            }
        }
    }
    return optimalPoint;
}
\end{verbatim}

\textbf{Hướng dẫn cài đặt:}
\begin{itemize}[leftmargin=*,noitemsep]
\item Không cần cài đặt. Chỉ cần gửi file \texttt{SmartLogi.html} cho học sinh.
\item Yêu cầu mở bằng trình duyệt web (Chrome, Firefox, Edge...).
\end{itemize}
\end{competitionsection}

\end{document}
